\clearpage
\section{Supervisor / Privileged Programming Model}

The privileged programming model defines how execution moves between user mode and supervisor mode, which saved
state is used for return, and which control state is directly accessible.

Supervisor entry changes the execution context at a single architectural commit point. If validation of the entry
target, saved state, or control-state write fails, the partial state transition is not exposed.

\subsection{Privileged Model Rules}
Privileged state is visible through STATUS, segment registers, control registers, interrupt state, and
address-translation controls. User-mode instructions may read or write only the state explicitly marked
unprivileged. Privileged writes validate reserved bits, selector values, and control-state consistency before
commit.

STATUS.PM is zero in user mode and one in supervisor mode. SYSCALL, trap, interrupt, and exception entry set it to
one at the entry commit point. SYSRET and IRET restore a saved STATUS image only after validating that the requested
return state is architecturally permitted. SYSRET is reserved for the SYSCALL frame shape; interrupt, trap, and
exception frames return through IRET.

SYSCALL saves FRAME\_CONTROL, SS:SP, CS:PC, and DS before loading the supervisor target from SPC, SCS, and SDS.
RDCR, WRCR, translation-state changes, interrupt-control changes, and page-table context switches require supervisor
privilege unless an instruction explicitly says otherwise. Delivery of a reserved CPU-owned vector is invalid and
enters the double-fault path.

\subsection{Privilege State}
STATUS.PM selects the current privilege level. In user mode, current-domain and user-domain memory accesses are
identical. In supervisor mode, ordinary memory operands use current-domain access. Checked user-domain access is
expressed only by the supervisor-only \texttt{MOVUC}, \texttt{MOVCU}, and \texttt{MOVUU} instructions.

\begin{manualtabular}{@{}>{\raggedright\arraybackslash}p{1.05in}>{\raggedright\arraybackslash}p{0.45in}>{\raggedright\arraybackslash}p{3.90in}@{}}
\toprule
\textbf{Field} & \textbf{Value} & \textbf{Meaning}\\
\midrule
STATUS.PM & 0 & user mode\\
STATUS.PM & 1 & supervisor mode\\
STATUS.IE & 0/1 & maskable interrupt enable\\
STATUS.TF & 0/1 & single-step trace control\\
STATUS.RF & 0/1 & resume flag for exception restart\\
STATUS.NI & 0/1 & nested-interrupt inhibit\\
STATUS.IN & 0/1 & inside interrupt or supervisor entry\\
current memory domain & -- & default domain for ordinary memory operands\\
checked user access & -- & MOVUC, MOVCU, and MOVUU select user-domain memory access\\
\bottomrule
\end{manualtabular}
\manualunlistedtablecaption{Privilege State and Access Domains}

\subsection{SYSCALL and SYSRET}
SYSCALL is an explicit supervisor-entry path rather than an IVT-vector event. It saves the user control state,
loads the supervisor entry state from SPC/SCS/SDS, and sets STATUS.PM to supervisor mode.

SYSRET is the matching return instruction for the syscall frame shape. It restores the saved user execution
context only after the saved control state has passed return validation. CS and PC changes are a single
architectural commit point for long control-transfer and supervisor-entry operations.

SYSCALL does not allocate or consume an IVT vector. Its saved frame contains FRAME\_CONTROL, SS:SP, CS:PC, and DS,
and the entry commit sets STATUS.PM to one. SYSRET reads that exact frame, validates the saved status and segment
state, and then restores the saved PC and execution context as one architectural transition.

@SUPERVISOR_CONTROL_FLOW_TABLE@

\subsection{Interrupt and Control-State Rules}
Interrupt delivery saves the interrupted PC, CS, FLAGS, and STATUS state needed for precise return. The entry
path validates the target vector and control state before making the supervisor state visible.

Writes to control state that would create an invalid translation, privilege, or interrupt configuration raise the
architecturally named fault without committing a partial control-register update.

The SN field in each IVT entry selects the supervisor stack set used to construct the entry frame; NMI and
double-fault delivery use the same mechanism. When the hidden current interrupt depth reaches ICR.MAX\_IDEPTH,
maskable interrupts are implicitly masked. Reserved control and status bits must validate before commit, and an
invalid write raises INVALID\_CONTROL\_STATE. User-readable control state is exposed through dedicated gated
instructions rather than direct control-register access.
